\documentclass[12pt]{book}
\usepackage{amsmath}
\usepackage{amssymb}
\usepackage{geometry}
\geometry{margin=1in}

\title{Volume 04: Electromagnetism and Magnetic Phenomena \\ Book 01: Electricity from Matrix Pressure}
\author{James C. Harvey}
\date{December 2025}

\begin{document}
\maketitle
\tableofcontents

\chapter{Electricity as Pressure Geometry}

\section*{Abstract}
This book derives electric phenomena as consequences of pressure geometry in the matrix. Electric field,
potential, and charge are geometric proxies for occlusion and curvature in the small-boundary regime.
Equations are expressed in SDT-native form with explicit constants.

\section{Electric Field from Pressure Gradient}
The electric field is the local pressure gradient expressed in force per unit proxy charge:
\begin{equation}
\mathbf{E} = -\nabla \Phi,
\end{equation}
with potential defined by
\begin{equation}
\Phi(\mathbf{r}) = -\int_{\infty}^{\mathbf{r}} \mathbf{E}\cdot d\mathbf{l}.
\end{equation}
In SDT, $\Phi$ is a pressure-normalized potential and therefore a flux geometry measure.

\section{Charge as Geometric Proxy}
Charge is not primitive. It is the small-boundary proxy for curvature:
\begin{equation}
\frac{\kappa^2}{4\pi K_{\text{bulk}}} = k_e q^2.
\end{equation}
This mapping recovers Coulomb’s law without introducing charge as a fundamental substance.

\section{Proxy Conversion}
Solving for $q$ yields the explicit conversion:
\begin{equation}
q = \frac{\kappa}{\sqrt{4\pi K_{\text{bulk}} k_e}}.
\end{equation}
This relation is the bridge between curvature geometry and charge-like bookkeeping.

\section{Coulomb Regime}
For occlusion $E \to 0$,
\begin{equation}
F_{\text{C}} = \frac{\kappa_1 \kappa_2}{4\pi K_{\text{bulk}} r^2}.
\end{equation}
Thus electric interaction is a regime of the same flux-driven coupling used everywhere in SDT.

\chapter{Electric Potential and Work}

\section*{Abstract}
Electric potential is the work per unit curvature proxy required to move a boundary in the matrix. This
chapter formalizes potential as pressure geometry and relates it to energy transfer.

\section{Potential as Pressure Integral}
\begin{equation}
\Phi(\mathbf{r}) = \frac{1}{\kappa_{\text{ref}}}\int_{\infty}^{\mathbf{r}} \nabla P \cdot d\mathbf{l}.
\end{equation}
The reference curvature $\kappa_{\text{ref}}$ sets the proxy conversion and is fixed by geometry.

\section{Work in the Matrix}
The work to move a boundary element through the matrix is
\begin{equation}
W = \int \mathbf{F}\cdot d\mathbf{l} = \kappa_{\text{ref}} \Delta \Phi.
\end{equation}

\chapter{Flux Continuity and Gauss Structure}

\section*{Abstract}
Gauss-like relations arise from flux conservation. Charge is a proxy for curvature, so enclosed curvature
determines flux divergence.

\section{Divergence Relation}
Define effective enclosed curvature $\kappa_{\text{enc}}$:
\begin{equation}
\nabla \cdot \mathbf{E} = \frac{\kappa_{\text{enc}}}{\epsilon_{\text{eff}}},
\end{equation}
where $\epsilon_{\text{eff}}$ encodes matrix compliance.

\chapter{Current as Flux Drift}

\section*{Abstract}
Current is flux drift along a boundary lattice. It is the macroscopic transport of matrix momentum
through constrained channels.

\section{Current Density}
\begin{equation}
\mathbf{J} = \rho_{\kappa} \mathbf{v}_d,
\end{equation}
where $\rho_{\kappa}$ is curvature proxy density and $\mathbf{v}_d$ is drift velocity along the boundary.

\section{Ohmic Behavior}
Ohmic behavior is the linear regime of drift under pressure gradient:
\begin{equation}
\mathbf{J} = \sigma \mathbf{E},
\end{equation}
with conductivity $\sigma$ determined by boundary smoothness and locking efficiency.

\section{Conductivity Scaling}
An SDT scaling for conductivity is
\begin{equation}
\sigma \sim \frac{\rho_{\kappa} \ell_c}{\lambda},
\end{equation}
where $\ell_c$ is a characteristic contact length. Higher locking efficiency lowers conductivity.

\chapter{Dielectric Response}

\section*{Abstract}
Dielectrics are media where boundaries polarize without sustained drift. This polarization is a local
rearrangement of occlusion and curvature.

\section{Polarization}
\begin{equation}
\mathbf{P} = \chi_e \mathbf{E},
\end{equation}
where $\chi_e$ is the SDT susceptibility defined by boundary response to pressure gradients.

\section{Energy Density}
The energy density stored in polarization is
\begin{equation}
u = \frac{1}{2}\mathbf{E}\cdot\mathbf{P} = \frac{1}{2}\chi_e E^2,
\end{equation}
which is the geometric work required to reconfigure boundary orientations.
\end{document}
